% Part: turing-machines
% Chapter: machines-computations
% Section: combining-machines

\documentclass[../../../include/open-logic-section]{subfiles}

\begin{document}

\olfileid{tur}{mac}{cmb}
\olsection{টুরিং যন্ত্রের সংযোজন}

\begin{explain}
এ পর্যন্ত দেখা টুরিং যন্ত্রের উদাহরণগুলি বেশ সরল। কিন্তু বাস্তবে আধুনিক
যেকোনো প্রোগ্রামিং ভাষায় সমাধানযোগ্য প্রত্যেক সমস্যা টুরিং যন্ত্রেও
সমাধান করা যায়। আরও জটিল টুরিং যন্ত্র তৈরি করতে যন্ত্রগুলি যে সংযোজন করা
যায়, সে বিষয়ে নিশ্চিত হওয়া জরুরি; তাহলে প্রক্রিয়াটিকে সহজতর অংশে ভেঙে
জটিলতর সমস্যার সমাধানকারী যন্ত্র তৈরি করা যাবে। কোনো জটিল সমস্যাকে তার
স্বাভাবিক উপাংশে ভাঙার উপায় পেলে সমস্যাটি কয়েকটি পর্যায়ে সামলানো যায়:
কয়েকটি সরল টুরিং যন্ত্র তৈরি করে সেগুলিকে একটি সমস্যাসমাধানকারী যন্ত্রে
সংযোজন করা যায়। পরের অংশে থামা সমস্যা মোকাবিলা করার সময় এই কথাটি
বিশেষভাবে গুরুত্বপূর্ণ।

টুরিং যন্ত্র $M = \tuple{Q, \Sigma, q_0, \delta}$ এবং
~$M' = \tuple{Q', \Sigma', q_0', \delta'}$-কে কীভাবে সংযোজন করব? এখন
$M$ থামার পর ফিতার কনফিগারেশনটিকে যন্ত্র~$M'$-এর একটি চালনের নিবেশ-
কনফিগারেশন হিসেবে ব্যবহার করি। এই কাজ করে এমন একটিমাত্র টুরিং যন্ত্র
$M \frown M'$ পেতে নিচের পদক্ষেপগুলি নাও:
\begin{enumerate}
    \item $M'$-এর সব দশা~$Q'$-কে নতুন সংখ্যা (বা নতুন নাম) দাও, যাতে
    $M$ ও~$M'$-এর কোনো দশা অভিন্ন না হয় ($Q \cap Q' = \emptyset$)।
    \item $M \frown M'$-এর দশা-সেট হলো $Q \cup Q'$।
    \item ফিতার বর্ণমালা হলো $\Sigma \cup \Sigma'$।
    \item আরম্ভিক দশা হলো~$q_0$।
    \item অবস্থান্তর অপেক্ষক হলো নিচে দেওয়া অপেক্ষক~$\delta''$:
    \[\delta''(q,\sigma) =
    \begin{cases}
      \delta(q,\sigma) & \text{যদি $q \in Q$ এবং $\delta(q,\sigma)$ সংজ্ঞায়িত হয়}\\
      \delta'(q,\sigma) & \text{যদি $q \in Q'$}\\
      \tuple{q_0', \sigma, \TMstay} & \text{যদি $q \in Q$ এবং
      $\delta(q,\sigma)$ অসংজ্ঞায়িত হয়}
    \end{cases}\]
\end{enumerate}
% BN-SRC-171 TARGET-CORRECTION-BEGIN
\par\smallskip\noindent\textnormal{\small\textbf{উৎস-সংশোধন (BN-SRC-171):} হিমায়িত ইংরেজি খণ্ডিত সংজ্ঞার প্রথম শর্তটি শুধু $q\in Q$ বলে, অথচ তৃতীয় শর্তটি $q\in Q$ এবং $\delta(q,\sigma)$ অসংজ্ঞায়িত বলে; ফলে তৃতীয় ক্ষেত্রে দুটি শর্ত একসঙ্গে প্রযোজ্য হয় এবং $\delta''$ অপেক্ষক হিসেবে সুসংজ্ঞায়িত থাকে না। বর্ণিত যন্ত্র-সংযোজন অনুযায়ী বাংলা পাঠে প্রথম শর্তে $\delta(q,\sigma)$ সংজ্ঞায়িত হওয়ার শর্ত যোগ করা হয়েছে।} % BN-SRC-171 TARGET-CORRECTION-END
ফলে পাওয়া যন্ত্রটি দশা $q \in Q$-তে থাকলে~$M$-এর নির্দেশ এবং দশা
$q \in Q'$-তে থাকলে~$M'$-এর নির্দেশ ব্যবহার করে। যন্ত্রটি যদি দশা
$q \in Q$-তে থেকে এমন প্রতীক~$\sigma$ পড়ে যার জন্য~$M$-এর কোনো
অবস্থান্তর নেই (অর্থাৎ $M$ থেমে যেত), তবে সেটি~$M'$-এর আরম্ভিক দশায়
প্রবেশ করে (আর ফিতার বিষয়বস্তু ও হেডের অবস্থান অপরিবর্তিত রাখে)।

খেয়াল করো, যন্ত্র~$M$ শৃঙ্খলাবদ্ধ না হলে $M$ থামার সময় ফিতার হেডটি
কোথায় থাকে তা আমরা জানি না; তাই~$M$-এর থামা-কনফিগারেশনে হেডটি ঘর~$1$
পড়বে, এমন কথা নেই। যন্ত্র সংযোজনের সময় এটি মনে রাখা জরুরি।
\end{explain}

\begin{ex}
\emph{যন্ত্রের সংযোজন:} এমন একটি যন্ত্রের নকশা করব, যা দৈর্ঘ্য $n$ ও~$m$-
এর দুটি $\TMstroke$-খণ্ড নিবেশ হিসেবে নিয়ে শুরু করে এবং ফিতায়
$2(m+n)$টি $\TMstroke$-এর একটিমাত্র খণ্ড রেখে থামে। যন্ত্রটি তৈরি করতে
আমাদের পরিচিত দুটি যন্ত্র—যোগ-যন্ত্র ও দ্বিগুণকারী—সংযোজন করা যায়। প্রথমে
যোগ-যন্ত্রের দশা-চিত্র আঁকি।
\[
\begin{tikzpicture}[->,>=stealth',shorten >=1pt,auto,node distance=2.8cm,
                    semithick]
  \tikzstyle{every state}=[fill=none,draw=black,text=black]

  \node[initial,state] (A)              {$q_0$};
  \node[state]         (B) [right of=A] {$q_1$};
  \node[state]         (C) [right of=B] {$q_2$};

  \path (A) edge node {\TMtrans{\TMblank}{\TMstroke}{\TMstay}} (B)
            edge [loop above] node {\TMtrans{\TMstroke}{\TMstroke}{\TMright}} (A)
        (B) edge [loop above] node {\TMtrans{\TMstroke}{\TMstroke}{\TMright}} (B)
            edge node {\TMtrans{\TMblank}{\TMblank}{\TMleft}} (C)
        (C) edge [loop above] node {\TMtrans{\TMstroke}{\TMblank}{\TMstay}} (C);
\end{tikzpicture}
\]
দশা~$q_2$-তে থামার বদলে নির্গমটিকে দ্বিগুণ করার কাজ চালিয়ে যেতে চাই।
মনে করো, দ্বিগুণকারী যন্ত্রটি নিবেশের প্রথম দাগটি মুছে আলাদা নির্গমে দুটি
দাগ লেখে। ফিতার হেড যাতে যোগ-যন্ত্রের নির্গমের প্রথম দাগটি পড়ে, তা
নিশ্চিত করতে একটি নির্দেশ যোগ করি।
\[
\begin{tikzpicture}[->,>=stealth',shorten >=1pt,auto,node distance=2.8cm,
                    semithick]
  \tikzstyle{every state}=[fill=none,draw=black,text=black]

  \node[initial,state] (A)              {$q_0$};
  \node[state]         (B) [right of=A] {$q_1$};
  \node[state]         (C) [right of=B] {$q_2$};
  \node[state]         (D) [below left of=C] {$q_3$};
  \node[state]         (E) [below left of=D] {$q_4$};

  \path (A) edge node {\TMtrans{\TMblank}{\TMstroke}{\TMstay}} (B)
            edge [loop above] node {\TMtrans{\TMstroke}{\TMstroke}{\TMright}} (A)
        (B) edge [loop above] node {\TMtrans{\TMstroke}{\TMstroke}{\TMright}} (B)
            edge node {\TMtrans{\TMblank}{\TMblank}{\TMleft}} (C)
        (C) edge node {\TMtrans{\TMstroke}{\TMblank}{\TMleft}} (D)
        (D) edge [loop left] node {\TMtrans{\TMstroke}{\TMstroke}{\TMleft}} (D)
            edge node[left, xshift=-2mm] {\TMtrans{\TMendtape}{\TMendtape}{\TMright}} (E);
\end{tikzpicture}
\]
এবার নিবেশটি দ্বিগুণ করা সহজ—দ্বিগুণকারী যন্ত্রটিকে শুধু দশা~$q_4$-এর
সঙ্গে জুড়ে দিতে হবে। তার জন্য দ্বিগুণকারী যন্ত্রের দশাগুলির নতুন নাম দিতে
হবে, যাতে তাদের সংখ্যা~$q_0$-র বদলে~$q_4$ থেকে শুরু হয়—এতে দুটি আরম্ভিক
দশা তৈরি হবে না। চূড়ান্ত চিত্রটি \olref{fig:combined}-এর মতো হবে।
\begin{figure}
\[
\begin{tikzpicture}[->,>=stealth',shorten >=1pt,auto,node distance=2.8cm,
                    semithick]
  \tikzstyle{every state}=[fill=none,draw=black,text=black]
  \node[initial,state] (A)              {$q_0$};
  \node[state]         (B) [right of=A] {$q_1$};
  \node[state]         (C) [right of=B] {$q_2$};
  \node[state]         (D) [below left of=C] {$q_3$};
  \node[state]         (E) [below left of=D] {$q_4$};
  \node[state]         (2) [right of=E] {$q_5$};
  \node[state]         (3) [right of=2] {$q_6$};
  \node[state]         (4) [below of=3] {$q_7$};
  \node[state]         (5) [left of=4]  {$q_8$};
  \node[state]         (6) [left of=5]  {$q_9$};

  \path (A) edge node {\TMtrans{\TMblank}{\TMstroke}{\TMstay}} (B)
            edge [loop above] node {\TMtrans{\TMstroke}{\TMstroke}{\TMright}} (A)
        (B) edge [loop above] node {\TMtrans{\TMstroke}{\TMstroke}{\TMright}} (B)
            edge node {\TMtrans{\TMblank}{\TMblank}{\TMleft}} (C)
        (C) edge node {\TMtrans{\TMstroke}{\TMblank}{\TMleft}} (D)
        (D) edge [loop left] node {\TMtrans{\TMstroke}{\TMstroke}{\TMleft}} (D)
            edge node[left, xshift=-2mm] {\TMtrans{\TMendtape}{\TMendtape}{\TMright}} (E)
    (E) edge node {\TMtrans{\TMstroke}{\TMblank}{\TMright}} (2)
    (2) edge [loop above] node {\TMtrans{\TMstroke}{\TMstroke}{\TMright}} (2)
      edge node {\TMtrans{\TMblank}{\TMblank}{\TMright}} (3)
    (3) edge [loop above] node {\TMtrans{\TMstroke}{\TMstroke}{\TMright}} (3)
        edge node {\TMtrans{\TMblank}{\TMstroke}{\TMright}} (4)
    (4) edge [loop below] node {\TMtrans{\TMblank}{\TMstroke}{\TMleft}} (4)
        edge node {\TMtrans{\TMstroke}{\TMstroke}{\TMleft}} (5)
    (5) edge [loop below]  node {\TMtrans{\TMstroke}{\TMstroke}{\TMleft}} (5)
        edge              node {\TMtrans{\TMblank}{\TMblank}{\TMleft}} (6)
    (6) edge [loop below] node {\TMtrans{\TMstroke}{\TMstroke}{\TMleft}} (6)
        edge              node {\TMtrans{\TMblank}{\TMblank}{\TMright}} (E);
\end{tikzpicture}
\]
\caption{যোগ-যন্ত্র ও দ্বিগুণকারী যন্ত্রের সংযোজন}
\ollabel{fig:combined}
\end{figure}
% BN-SRC-172 TARGET-CORRECTION-BEGIN
\par\smallskip\noindent\textnormal{\small\textbf{উৎস-সংশোধন (BN-SRC-172):} এই উদাহরণের হিমায়িত ইংরেজি পাঠে যোগ-যন্ত্রের তিনটি পুনরাবৃত্ত চিত্রেই $q_0$-তে $\TMstroke$ পড়ার তীরটি ভুল করে $q_1$-এ যায়। বর্ণিত যোগ-পদ্ধতি অনুযায়ী বাংলা চিত্রগুলির তিনটিতেই তীরটি $q_0$-র স্ব-আবর্তন করা হয়েছে।} % BN-SRC-172 TARGET-CORRECTION-END
\end{ex}

\begin{prop}
    যদি $M$ ও~$M'$ শৃঙ্খলাবদ্ধ হয় এবং যথাক্রমে অপেক্ষক
    $f\colon \Nat^k \to \Nat$ ও $f'\colon \Nat \to \Nat$ গণনা করে,
    তবে $M \frown M'$ শৃঙ্খলাবদ্ধ এবং~$\comp{f}{f'}$ গণনা করে।
\end{prop}

\begin{proof}
    $M$ শৃঙ্খলাবদ্ধ বলে সেটি $f(n_1,\dots,n_k) = m$ নির্গম নিয়ে থামার
    সময় হেডটি ঘর~$1$ পড়ে। এবার~$M'$-এর আরম্ভিক দশায় প্রবেশ করলে যন্ত্রটি
    আবার ঘর~$1$ পড়তে পড়তে নির্গম $f'(m)$ নিয়ে থামবে।
    \olref[dis]{defn:disciplined}-এর অন্য শর্তগুলিও পূরণ হয়।
\end{proof}

\begin{prob}
    \cref{tur:mac:dis:prob:disc-succ}-এর যন্ত্র~$M$ নিয়ে $M \frown M$
    নির্মাণ করে $f(x) = x+2$ গণনাকারী একটি শৃঙ্খলাবদ্ধ টুরিং যন্ত্র দাও।
\end{prob}
\end{document}
